% Shared math macros for MATH 223 (Linear Algebra).
% Loaded by BOTH slides and notes (via slides-style.tex / notes-style.tex),
% so notation always matches between the two.
% Notation follows the existing decks (e.g. Set 7).

\usepackage{amsmath,amssymb,mathtools}

% ---------- Number systems & spaces ----------
\newcommand{\R}{\mathbb{R}}
\newcommand{\C}{\mathbb{C}}
\newcommand{\Z}{\mathbb{Z}}
\newcommand{\Q}{\mathbb{Q}}
\newcommand{\Mat}[2]{M_{#1\times #2}}          % \Mat{2}{2}  -> M_{2x2}
\newcommand{\Poly}[1]{P_{#1}}                  % \Poly{2}    -> P_2

% ---------- Vectors ----------
\newcommand{\vect}[1]{\mathbf{#1}}             % \vect{v}
\newcommand{\zero}{\mathbf{0}}
\newcommand{\e}[1]{\mathbf{e}_{#1}}            % \e{1}
\newcommand{\colvec}[1]{\begin{bmatrix} #1 \end{bmatrix}}  % \colvec{1\\2\\3}
\newcommand{\T}{^{\mathsf{T}}}                 % A\T
\newcommand{\inv}{^{-1}}
\newcommand{\norm}[1]{\left\lVert #1 \right\rVert}
\newcommand{\ip}[2]{\left\langle #1,\, #2 \right\rangle}
\newcommand{\abs}[1]{\left\lvert #1 \right\rvert}

% ---------- Bases, coordinates, matrices of maps ----------
\newcommand{\B}{\mathcal{B}}
\newcommand{\Cb}{\mathcal{C}}                  % (\C is the complex numbers)
\newcommand{\E}{\mathcal{E}}
\newcommand{\coord}[2]{\bigl[#1\bigr]_{#2}}    % \coord{v}{\B}        -> [v]_B
\newcommand{\matof}[3]{\bigl[#1\bigr]_{#2\to #3}} % \matof{T}{\B}{\Cb} -> [T]_{B->C}
\newcommand{\Dmap}[1]{D_{#1}}                  % deconcretization map D_B

% ---------- Operators ----------
\DeclareMathOperator{\im}{im}
\DeclareMathOperator{\col}{col}
\DeclareMathOperator{\Row}{row}
\DeclareMathOperator{\Span}{span}
\DeclareMathOperator{\rank}{rank}
\DeclareMathOperator{\nullity}{nullity}
\DeclareMathOperator{\tr}{tr}
\DeclareMathOperator{\proj}{proj}
\DeclareMathOperator{\RREF}{RREF}
\newcommand{\charpoly}[1]{C_{#1}}              % \charpoly{T}(x)
\newcommand{\Eig}[1]{E_{#1}}                   % \Eig{\lambda}
\newcommand{\mg}{m_g}
\newcommand{\ma}{m_a}

% ---------- Row reduction ----------
% Augmented matrix: \begin{amatrix}{3} 1 & 2 & 3 & 4 \\ ... \end{amatrix}
% (argument = number of coefficient columns)
\newenvironment{amatrix}[1]{%
  \left[\begin{array}{@{}*{#1}{c}|c@{}}
}{%
  \end{array}\right]
}
\newcommand{\rowop}[1]{\xrightarrow{\;#1\;}}   % \rowop{R_3-2R_1}
